% ============================================================
% 1 引言
% ============================================================
\section{引言}

\subsection{从“检测”到“溯源”}

深度伪造技术——换脸、表情重演与语音合成——已对个人隐私、司法取证与舆情治理构成实质威胁
\cite{rossler2019faceforensics}。过去数年，主流研究围绕\textbf{真假判别}展开，
即把图像或视频二分类为真或假。然而在真实取证流程中，“这是假的”往往只是起点：
调查者真正需要回答的是\emph{“这张伪造图沿用了哪一张真实图的背景？”}
——从生成内容反查其\textbf{真实来源}。

一个可行的事实基础是：主流换脸算法通常\textbf{保留}目标视频的背景与头姿，只替换前景人脸，
因此\textbf{背景}构成了连通伪造图与其真实源图的\emph{来源指纹}。
若在传播发生之前先对原始真实素材建立索引，之后便可由一张伪造图检索回它的来源。
这一思路被称为\emph{主动}（proactive）溯源，与事后的被动检测互补。
本文沿用此设定，称该任务为 \aura{}（Attack-Unsupervised Retrieval of Authentic backgrounds）；
“\emph{无攻击监督}”指训练阶段不使用任何篡改类型标签，以应对不断演化的未知篡改方式。

\subsection{已有的溯源工作}

Liu 等\cite{liu2026hashvae} 提出了图像级的 HashVAE：将待查图分为前景与背景分别编码，
以多尺度多分支残差主干学习紧凑哈希码，并用四元组损失对齐查询与库两侧的码空间。
Pei 等\cite{pei2021vithash} 则把问题做到视频级（ViTHash，即 HashViT）：
以视觉 Transformer 学习哈希中心，配合 Hash Triplet Loss 处理“同场景、相似人物”的难分样本，
并在检索到源视频后由定位模块给出篡改区域。

两项工作的共同设定是：库与查询都是\emph{干净}的、\emph{闭集}的、并且往往\emph{每视频只取一帧}；
评价指标以 Top-5 为主，检索方向也不完全一致（例如 HashVAE 以真实图查询伪造库）。
在这些设定下，学习式哈希已取得很高的读数。本文关注的是另一组更接近取证现场的问题：
当待查件被转存、压缩、裁剪过，当它的源头根本不在库中，此时方法的表现如何？
以及，在这样一个以背景低频信息为主的任务上，方法的收益究竟来自学习，还是来自描述子本身？

\subsection{本文的问题、假设与验证路线}

\chead{问题。}在该任务上，视频级与帧级呈现出截然不同的难度：只用零训练描述子时，
视频级已达 $99\%$ 以上，而帧级只有 $60\%$ 左右；我们训练的学习编码器最好也只到 $60.05\%$，
与零训练基线持平，并未超过。两个层级的差别在于：视频级只需识别“背景属于哪个视频”，
帧级却要在同一视频的 8 张背景几乎逐像素相同的帧中挑出正确的那一张。
换句话说，帧级所需的判别信号是\textbf{微小而局部}的。

\chead{假设。}我们假设帧级失效并非容量不足，而是建模过程主动抹去了这种微小信号。
具体有三条途径：（a）嵌入头之前的\textbf{全局池化}把空间布局压成一个向量；
（b）训练时的\textbf{几何增强}显式要求模型对旋转、缩放、平移不变，
而帧级恰恰要靠这些空间位置来区分；（c）以\textbf{单一全局向量}做比对，
一次只能给出“整体像不像”，无法校验局部是否对得上。

\chead{证据。}四条读数支持该假设：其一，去掉全局池化后帧级由约 $46\%$ 跃升至约 $58\%$
（$+12$ 个百分点），说明有用信息确实在空间布局里；其二，加入几何增强的一轮反而下降，
说明“对几何不变”与“靠几何区分”是相反的训练方向；其三，
完全不训练、但保留 $8\times8$ 空间网格并逐图标准化的描述子达到 $60.05\%$，
说明\textbf{保留空间网格本身}就是关键，而与是否学习无关；其四，
把候选集限制为查询自身源视频的 $8$ 帧后，学习编码器与零训练描述子命中率持平，
即帧级的全部误差都在同视频内部，学习并未带来额外分辨力。
进一步的单变量训练排除了“负样本不够难”这一解释：把同视频相邻帧固定放进批内，读数并未提升。
第二轮单变量训练在全局向量之外补了一条逐位置监督（$8\times8$ 网格逐位置对比），
$8$ 选 $1$ 读数为 $60.14\%$，与零训练描述子的 $60.11\%$ 持平、差值在轮次波动之内；
由于检索仍以单一全局内积排序，这一结果说明“编码器看不到细节”并不成立，
瓶颈更可能在\textbf{检索时的比较方式}。零训练分辨率扫描也与此一致：
$8\times8$ 即为峰值，$16\times16$ 以上反而递减，只保留高频的输入仅 $46.74\%$，
因此“更细的细节”并不是缺失的那部分信息。

\chead{从任务性质到方法。}上述诊断给出了方法必须满足的约束，而不是一个待调的架构：
背景既是唯一可用信号，就不能只保留整图的高层语义（P1）；
判别信息在低频布局而非高频纹理，就不能靠加深网络或提高分辨率换取（P2）；
帧级要靠局部位置差异区分，就不能用单一全局内积完成比较（P3）；
几何退化会破坏对应关系，就需要显式的成对几何校验（P4）；
而在光度退化下这类校验没有可补的信息，因此需要按条件启用而非无条件叠加（P5）；
库中可能根本没有来源，检索与拒识就必须在同一个相似度上联合判断（P6）。
把六条约束逐条落成模块，即第 4 节的方法：背景加权、低频友好的全局编码器、
可微的局部网格成对打分、带门控与几何回退的两阶段重排，以及基于同一 top-1 相似度的开集拒识。

\chead{验证与结果。}按这一设计，在标准 $8$ 帧协议下，干净档帧级命中由初筛的 $60.00\%$ 提升至 $67.67\%$
（$+7.7$ 个百分点），几何七档均值 $66.90\%$，与经典 SIFT$+$RANSAC 的 $67.57\%$ 相当。
更重要的是，我们按 P2 与 P5 的预期检验了各模块：在微小差异本来就显式的稠密帧库上，
全局外观已足以定位，几何重排的增益只出现在几何退化档（第 5.7 节）——
这与“几何校验补的是对应关系、而非分辨力”的预期一致，也说明方法的价值应按条件分档陈述。

\chead{协议。}上述结论需要一套统一的评测条件才可信，因此我们同时固定了
图库$/$查询的构造与帧级$/$视频级判据，并补齐两类真实条件：
只作用于查询的\textbf{确定性退化}（光度六档 $+$ 几何六档），
以及“库中可能无源”的\textbf{开集拒识}。

这条“问题 $\rightarrow$ 假设 $\rightarrow$ 证据 $\rightarrow$ 手段 $\rightarrow$ 结果”的链条
即本文的主线；第 5.3 节陈述假设与证据，第 5.6 节给出验证结果。

\subsection{主要贡献}

\begin{enumerate}[leftmargin=1.6em,itemsep=2pt]
\item \textbf{任务与可复现协议。} 形式化 \aura{}（无攻击监督的真实源图检索），
      给出图库$/$查询构造、帧级与视频级命中判据及随机基线，并实现为可复现代码（第 3 节）。
\item \textbf{退化鲁棒基准。} 一族只作用于查询、图库保持干净的确定性退化协议
      （光度六档 $+$ 几何六档），并在其上得到明确结论：
      光度退化几乎不改变读数，几何退化才是该任务的有效评测条件（第 3.4、5.4 节）。
\item \textbf{真帧级定位协议。} 现有“帧级”判据实为从同一视频抽出的 $8$ 帧中挑一张。
      我们构造稠密帧库（$30$ 条视频各连续 $100$ 帧，逐帧检测人脸），查询伪造帧，
      并把命中范围定义为\emph{相似度区间}而非固定帧号，以吸收换脸方法的时间轴漂移
      与静止段的帧间歧义；在该协议下重新界定“帧级”可达到的精度（第 5.7 节）。
\item \textbf{开集拒识协议。} 引入闭集$/$开集切分与以 top-1 相似度为判据的拒识 AUC$/$TPR@FPR；
      实测显示拒识与检索呈现同一模式——对光度退化近乎不变，对几何退化显著下降（第 4.5、5.5 节）。
\item \textbf{帧级失效的成因假设与验证。} 提出“建模把微小空间信号交给单一全局向量的内积比对”的假设，
      以四组读数支持之（去池化 $+12$ 个百分点、几何增强反向、零训练空间网格最强、
      同视频 $8$ 选 $1$ 与零训练持平），并以两组单变量训练排除“负样本不够难”的解释、
      以及“编码器看不到细节”的解释（逐位置监督只带来轮次波动量级的差别）；
      在第 5.3、5.6 节完成从假设到验证的闭环：把整图比对换成细粒度局部网格上的成对比较后，
      干净档帧级由 $60.00\%$ 升至 $67.67\%$。
\item \textbf{可微几何验证路径。} 给出一条由软互最近邻、可微 DLT 与软内点占比构成的
      两阶段重排器，在几何七档上与经典 SIFT$+$RANSAC 并列比较，二者表现相当，
      并附参数敏感性分析（第 4.3、5.6、5.8 节）。
\end{enumerate}

\subsection{论文组织}

第 2 节回顾相关工作；第 3 节给出任务定义、数据划分与退化基准；
第 4 节描述编码器、可微几何重排器与开集拒识；第 5 节报告全部实验读数与消融；
第 6 节讨论并总结。
